
\documentclass{article}
\usepackage{colortbl}
\usepackage{makecell}
\usepackage{multirow}
\usepackage{supertabular}

\begin{document}

\newcounter{utterance}

\centering \large Interaction Transcript for game `eqbench', experiment `eq\_bench', episode 166 with reflexion{-}Qwen3.5{-}9B.
\vspace{24pt}

{ \footnotesize  \setcounter{utterance}{1}
\setlength{\tabcolsep}{0pt}
\begin{supertabular}{c@{$\;$}|p{.15\linewidth}@{}p{.15\linewidth}p{.15\linewidth}p{.15\linewidth}p{.15\linewidth}p{.15\linewidth}}
    \# & \multicolumn{2}{c}{Player} && \multicolumn{2}{c}{Game Master} \\
    \hline

    \theutterance \stepcounter{utterance}  
    & & & \multicolumn{4}{p{0.6\linewidth}}{
        \cellcolor[rgb]{0.9,0.9,0.9}{
            \makecell[{{p{\linewidth}}}]{
                \texttt{\tiny{[P1$\langle$GM]}}
                \\ 
\texttt{Your task is to predict the likely emotional responses of a character in this dialogue:} \\
\\ 
\texttt{Jonny: You're not like the others, are you, Eva?} \\
\texttt{Eva: And what makes you say that?} \\
\texttt{Jonny: Well, for starters, you can see me. And you don't scream and run away.} \\
\texttt{Eva: (laughs) I suppose that does make me a bit different.} \\
\texttt{Jonny: More than a bit, I'd say. } \\
\texttt{Eva: I've always been... sensitive to things others can't see. } \\
\texttt{Jonny: Like a sixth sense?} \\
\texttt{Eva: Something like that. } \\
\texttt{Jonny: Must be lonely, seeing things no one else can.} \\
\texttt{Eva: It can be... but it also allows me to meet interesting individuals. Like you.} \\
\texttt{Jonny: (laughs) I'm flattered. Not many people find a ghost interesting.} \\
\texttt{Eva: You're more than just a ghost, Jonny. You're a person, just... not in the traditional sense.} \\
\texttt{Jonny: That's... unexpectedly kind of you to say, Eva.} \\
\texttt{Eva: (smiles) Well, it's true.} \\
\texttt{Jonny: (pauses) Do you... Do you ever get scared? Dealing with... well, things like me?} \\
\texttt{Eva: Sometimes. But not with you. You're different.} \\
\texttt{Jonny: Different how?} \\
\texttt{Eva: You're... (pauses) You're not like the others.} \\
\texttt{Jonny: (laughs) Now who's sounding like a broken record?} \\
\texttt{Eva: (laughs) I suppose I walked into that one.} \\
\texttt{Jonny: (softly) Eva... } \\
\texttt{Eva: Yes, Jonny?} \\
\texttt{Jonny: Never mind. It's... it's nothing.} \\
\texttt{{[}End dialogue{]}} \\
\\ 
\texttt{At the end of this dialogue, Eva would feel...} \\
\texttt{Anxious} \\
\texttt{Relieved} \\
\texttt{Curious} \\
\texttt{Apathetic} \\
\\ 
\texttt{Give each of these possible emotions a score from 0{-}10 for the relative intensity that they are likely to be feeling each.} \\
\\ 
\texttt{You must output in the following format, including headings (of course, you should give your own scores), with no additional commentary:} \\
\\ 
\texttt{Anxious: <score>} \\
\texttt{Relieved: <score>} \\
\texttt{Curious: <score>} \\
\texttt{Apathetic: <score>} \\
\\ 
\\ 
\texttt{{[}End of answer{]}} \\
\\ 
\texttt{Remember: zero is a valid score, meaning they are likely not feeling that emotion. You must score at least one emotion > 0.} \\
\\ 
\texttt{Your answer:} \\
            }
        }
    }
     \\ \\

    \theutterance \stepcounter{utterance}  
    & \multicolumn{4}{p{0.6\linewidth}}{
        \cellcolor[rgb]{0.561,0.737,0.561}{
            \makecell[{{p{\linewidth}}}]{
                \texttt{\tiny{[P1$\rangle$GM]}}
                \texttt{Anxious: 2} \\
\texttt{Relieved: 8} \\
\texttt{Curious: 6} \\
\texttt{Apathetic: 0} \\
            }
        }
    }
    & & \\ \\

    \theutterance \stepcounter{utterance}  
    & & & \multicolumn{2}{p{0.3\linewidth}}{
        \cellcolor[rgb]{0.9,0.9,0.9}{
            \makecell[{{p{\linewidth}}}]{
                \texttt{\tiny{[GM$|$GM]}}
                \texttt{Anxious: 2} \\
\texttt{Relieved: 8} \\
\texttt{Curious: 6} \\
\texttt{Apathetic: 0} \\
            }
        }
    }
    & & \\ \\

    \theutterance \stepcounter{utterance}  
    & & & \multicolumn{2}{p{0.3\linewidth}}{
        \cellcolor[rgb]{0.9,0.9,0.9}{
            \makecell[{{p{\linewidth}}}]{
                \texttt{\tiny{[GM$|$GM]}}
                \texttt{Anxious: 2} \\
\texttt{Relieved: 0} \\
\texttt{Curious: 8} \\
\texttt{Apathetic: 0} \\
            }
        }
    }
    & & \\ \\

    \theutterance \stepcounter{utterance}  
    & & & \multicolumn{2}{p{0.3\linewidth}}{
        \cellcolor[rgb]{0.9,0.9,0.9}{
            \makecell[{{p{\linewidth}}}]{
                \texttt{\tiny{[GM$|$GM]}}
                \texttt{game\_result = LOSE} \\
            }
        }
    }
    & & \\ \\

    \theutterance \stepcounter{utterance}  
    & & \multicolumn{4}{p{0.6\linewidth}}{
        \cellcolor[rgb]{0.95,0.95,0.95}{
            \makecell[{{p{\linewidth}}}]{
                \texttt{\tiny{[P1$\leftrightarrow$P]}}
                \texttt{Post{-}episode Reflexion (not sent to GM):} \\
\texttt{New plan: I failed because I assigned a high score to "Relieved" based on the character's polite acceptance, but I overlooked the underlying tension in Jonny's final "Never mind" and the heavy pause, which suggests unresolved fear or sadness rather than relief. Additionally, I gave "Curious" a moderate score despite the dialogue explicitly resolving the mystery of her nature, making curiosity irrelevant at the end. Next time, I will analyze the final emotional beat: Jonny's retreat ("Never mind") and Eva's gentle acceptance ("Yes, Jonny?") indicate a shift from fear to a quiet, shared understanding. I will lower "Relieved" to reflect that the moment is more poignant than happy, increase "Curious" only if there were lingering questions (which there aren't), and significantly boost "Anxious" to account for the lingering unease of the unspoken "things" Jonny mentioned, or better yet, recognize that the dominant feeling is a complex mix of **Relieved** (fear is gone) and **Curious** (about their future connection), but given the "LOSE" outcome, the most likely error was misidentifying the primary emotion. The correct strategy is to identify that the resolution of the "fear" question leads to a state of **Relieved** (high score) but tempered by the melancholy of the situation, while **Curious** should be low because the mystery is solved. However, since I scored Relieved 8 and Curious 6 and still lost, perhaps the intended answer prioritized the lingering **Anxious** feeling about the "things" or the awkwardness of the "Never mind." Let's re{-}evaluate: The prompt asks for the *likely* response. Eva just admitted she is different and that Jonny is different. The tension broke. "Relieved" is strong. "Curious" is weak. "Anxious" is low. "Apathetic" is 0. If I lost, maybe the "Curious" score was too high? Or perhaps the "Anxious" score should be higher due to the "things like me" context? Actually, looking at the dialogue, the "Curious" option is the most suspicious. The mystery is solved. The "Relieved" score of 8 seems appropriate. The "Anxious" score of 2 seems appropriate. The "Apathetic" score of 0 is correct. The only variable is if the system expects a different primary emotion. Wait, the instruction says "Your plan must change the strategy". My previous strategy was to guess based on surface politeness. The new strategy is to focus on the *resolution* of the conflict. The conflict was "Are you scared of ghosts?" {-}> "No, not with you." {-}> "Why?" {-}> "You're different." {-}> "Never mind." The emotional arc is Fear {-}> Relief {-}> Connection. The "Curious"} \\
            }
        }
    }
    & \\ \\

\end{supertabular}
}

\end{document}
